\chapter{Fungsi Pembangkit Peluang}\label{ch:b2:genfun}

Deret pangkat pada \cref{ch:b2:powerseries} kembali dengan misi
peluang: pada \omterm{def:b2:randomvar:law}{peubah acak} bernilai-$\N$ kita lekatkan
deret pangkat berkoefisien $\P(X = n)$. \emph{\omterm{ex:b2:powerseries:fibonacci}{Fungsi
pembangkit}} ini mengubah jumlah peubah yang \omterm{def:b2:proba:independence}{saling bebas}
menjadi hasil kali, momen menjadi turunan di $1$, dan kesamaan
kombinatorik yang sukar menjadi perkalian satu baris. Bab
ini menutup buku dengan dua pertunjukan: hampiran Poisson
bagi \omterm{def:b2:proba:space}{kejadian} langka, dan kriteria kepunahan bagi
\omterm{pb:b2:genfun:1}{proses percabangan} --- sebuah perhitungan peluang yang
sungguh tak hingga, yang terpecahkan seluruhnya oleh
geometri sebuah kurva \omterm{def:b2:affine:convex}{cembung}.

\section{Definisi dan sifat dasarnya}

\begin{definition}[Fungsi pembangkit peluang]\label{def:b2:genfun:pgf}
Misalkan $X$ \omterm{def:b2:randomvar:law}{peubah acak} bernilai-$\N$, dengan $p_n = \P(X = n)$. Adapun
\emph{fungsi pembangkit peluang}\index{fungsi pembangkit
peluang} bagi $X$ adalah jumlah
deret pangkatnya
\[
G_X(t) = \E\bigl(t^X\bigr) = \sum_{n=0}^{\infty} p_n\,t^n .
\]
\end{definition}

\begin{example}[Refleks pertama]\label{ex:b2:genfun:firstreflexes}
Peubah konstan $X = c$ punya $G_X(t) = t^c$; sebuah geseran menuruti
$G_{X+c}(t) = t^c\,G_X(t)$; dan mengevaluasinya di titik istimewa
membaca informasinya tanpa penguraian apa pun: $G_X(0) = \P(X
= 0)$, $G_X(1) = 1$, dan $G_X(-1) = \P(X\text{ genap}) -
\P(X\text{ ganjil})$, yakni neraca keparitasan yang dimanfaatkan pada
\cref{exo:b2:genfun:10}. Kalimat sebaris ini dipakai secara diam-diam
di mana-mana di bawah --- dan evaluasi $G_X(0)$ itu persis
cara peluang kepunahan kelak diperas dari \omterm{ex:b2:powerseries:fibonacci}{fungsi
pembangkit} yang teriterasi pada akhir bab ini.
\end{example}

\begin{proposition}[Jari-jari dan sifat pertamanya]\label{prop:b2:genfun:radius}
Deret yang mendefinisikan $G_X$ berjari-jari kekonvergenan $\geq 1$;
$G_X$ terdefinisi dan \omterm{def:b2:metric:continuity}{kontinu} pada $\intcc{-1}{1}$,
$\mathcal{C}^\infty$ pada $\intoo{-1}{1}$, dengan $G_X(1) = 1$ dan
$\abs{G_X(t)} \leq 1$ di sana. Lebih jauh, $G_X$ menentukan \omterm{def:b2:randomvar:law}{hukum}
$X$:
\[
p_n = \frac{G_X^{(n)}(0)}{n!} .
\]
\end{proposition}

\begin{proof}
Karena $\sum p_n = 1$ konvergen, sukunya $p_n\,1^n$ terbatas,
jadi jari-jarinya $\geq 1$ (lema Abel,
\cref{ch:b2:powerseries}); di $t = \pm1$ deretnya konvergen
\omterm{def:b2:series:def}{mutlak} ($\sum p_n = 1$ mendominasinya); lebih baik lagi, pada seluruh
selang $\intcc{-1}1$,
\[
\sup_{\abs t\leq1}\,\abs{p_nt^n} = p_n
\quad\text{dengan}\quad \sum_np_n < \infty :
\]
deretnya konvergen \emph{normal} pada $\intcc{-1}1$, jadi
jumlahnya \omterm{def:b2:metric:continuity}{kontinu} di sana
(\cref{thm:b2:funcseq:weierstrass,thm:b2:funcseq:continuity}). Adapun kemulusan di dalamnya dan rumus
koefisiennya merupakan teori umum deret pangkat; karena koefisiennya
dapat dipulihkan, dua peubah dengan \omterm{ex:b2:powerseries:fibonacci}{fungsi pembangkit} yang sama
berhukum sama.
\end{proof}

\begin{example}[Hukum klasiknya]\label{ex:b2:genfun:classical}
\leavevmode
\begin{itemize}
\item Bernoulli $\mathcal{B}(p)$: $G(t) = 1 - p + pt$.
\item Binomial $\mathcal{B}(n, p)$: $G(t) = \sum_k \binom nk
(pt)^k(1-p)^{n-k} = (1 - p + pt)^n$ (teorema binomial).
\item Geometrik $\mathcal{G}(p)$: $G(t) =
\sum_{k\geq1}(1-p)^{k-1}p\,t^k = \dfrac{pt}{1 - (1-p)t}$ (jari-jari
$\frac{1}{1-p} > 1$).
\item Poisson $\mathcal{P}(\lambda)$: $G(t) =
\sum_k e^{-\lambda}\frac{(\lambda t)^k}{k!} = e^{\lambda(t - 1)}$
(jari-jari $\infty$).
\end{itemize}
\end{example}

\begin{example}[Mengintegralkan fungsi
pembangkitnya]\label{ex:b2:genfun:integraltrick}
Turunan $G_X$ di $1$ memberi momen positif; adapun
\emph{integralnya} memberi yang negatif. Dari $\int_0^1t^k\dd t
= \frac1{k+1}$ dan pengintegralan suku demi suku (\omterm{def:b2:funcseq:series}{kekonvergenan
normal} pada $\intcc01$):
\[
\int_0^1G_X(t)\,\dd t = \sum_{k\geq0}\frac{\P(X =
k)}{k+1} = \E\Bigl(\frac1{1+X}\Bigr).
\]
Untuk $X \sim \mathcal P(\lambda)$:
\[
\E\Bigl(\frac1{1+X}\Bigr) =
\int_0^1\eu^{\lambda(t-1)}\,\dd t = \frac{1 -
\eu^{-\lambda}}{\lambda},
\]
yang memulihkan dalam satu baris perhitungan deret pada
\cref{ex:b2:randomvar:transferex}. \omterm{ex:b2:powerseries:fibonacci}{Fungsi pembangkitnya}
merupakan alat dua arah: turunkan di $1$ untuk memperoleh
momen $\E(X)$, $\E(X(X-1))$, integralkan pada $\intcc01$
untuk memperoleh $\E\bigl(\frac1{1+X}\bigr)$ --- satu objek \omterm{def:b2:powerseries:analytic}{analitik},
yang ditanyai ke arah mana pun yang diperlukan masalahnya.
\end{example}

\begin{example}[Sebuah hukum yang jari-jarinya persis
satu]\label{ex:b2:genfun:heavytail}
Ambil $\P(X = k) = \dfrac{6}{\pi^2k^2}$ untuk $k \geq 1$ --- sebuah
\omterm{def:b2:randomvar:law}{hukum} peluang menurut kesamaan Basel
(\cref{ex:b2:fourier:basel}). \omterm{ex:b2:powerseries:fibonacci}{Fungsi pembangkitnya}
$G(t) = \frac6{\pi^2}\sum_{k\geq1}\frac{t^k}{k^2}$ berjari-jari
kekonvergenan persis $1$: batas umum ``jari-jari $\geq
1$'' pada \cref{prop:b2:genfun:radius} tak dapat diperbaiki. Adapun
reratanya adalah
\[
\sum_{k\geq1}k\,\P(X = k) =
\frac6{\pi^2}\sum_{k\geq1}\frac1k = \infty :
\]
$G$ \omterm{def:b2:metric:continuity}{kontinu} pada $\intcc{-1}1$, mulus di dalamnya, tetapi
turunannya meledak di $1^-$ --- grafiknya tiba di
titik $(1, 1)$ dengan singgung tegak. Ekor yang berat terlihat
secara \emph{geometrik} pada \omterm{ex:b2:powerseries:fibonacci}{fungsi pembangkitnya}, di
satu titik $t = 1$ saja; adapun teorema momen di bawah membuat
padanan ini menjadi persis.
\end{example}

\begin{theorem}[Momen dari fungsi
pembangkitnya]\label{thm:b2:genfun:moments}
$X$ punya \omterm{def:b2:randomvar:expectation}{nilai harapan} bila dan hanya bila $G_X$ \omterm{def:b2:diffcalc:differential}{terdiferensialkan} di
$1^-$ (turunan kirinya, berhingga), dan lalu $\E(X) = G_X'(1)$.
Serupa itu $X$ punya momen kedua bila dan hanya bila $G_X$ dua kali
\omterm{def:b2:diffcalc:differential}{terdiferensialkan} di $1^-$, dan lalu
\[
\E\bigl(X(X - 1)\bigr) = G_X''(1),
\qquad
V(X) = G_X''(1) + G_X'(1) - G_X'(1)^2 .
\]
\end{theorem}

\begin{proof}
Untuk $t \in \intoo{0}{1}$, penurunan suku demi suku di dalam
cakramnya memberi $G_X'(t) = \sum_{n\geq1} np_n t^{n-1}$, yakni deret
berkoefisien taknegatif: $t \mapsto G_X'(t)$ tidak turun pada
$\intoo{0}{1}$, dan menurut kekonvergenan monoton jumlah parsialnya (atau
teorema Abel bagi koefisien taknegatif,
\cref{ch:b2:powerseries}),
\[
\lim_{t \to 1^-} G_X'(t)
= \sum_{n\geq1} n\,p_n \in \intcc{0}{+\infty} ,
\]
dengan setiap ruasnya berhingga persis ketika ruas satunya berhingga. Bila
berhingga, teorema nilai rata-rata mengapit hasil bagi selisihnya
$\frac{G_X(1) - G_X(t)}{1 - t}$ di antara nilai $G_X'$, jadi $G_X$
\omterm{def:b2:diffcalc:differential}{terdiferensialkan} di $1^-$ dengan $G_X'(1) = \sum np_n = \E(X)$ (menurut
pemindahannya). Pernyataan orde keduanya mengulangi hujahnya satu tingkat
lebih atas: $G''_X(t) = \sum_{n\geq2}n(n-1)p_nt^{n-2}$ juga
tidak turun pada $\intoo01$ dengan limit monoton
$\sum_nn(n-1)p_n = \E(X(X-1))$, yang berhingga persis ketika $X$ punya
momen kedua. Rumus \omterm{def:b2:randomvar:variance}{ragamnya} lalu menyusul dari
König--Huygens:
\[
V(X) = \E(X^2) - \E(X)^2 = \E\bigl(X(X-1)\bigr) + \E(X) -
\E(X)^2 = G''_X(1) + G'_X(1) - G'_X(1)^2 .
\]
\end{proof}

\begin{example}
Poisson: $G'(t) = \lambda e^{\lambda(t-1)}$, jadi $\E(X) = \lambda$;
$G''(1) = \lambda^2$, jadi $V(X) = \lambda^2 + \lambda - \lambda^2 =
\lambda$ --- yakni perhitungan pada \cref{ch:b2:randomvar} yang
masing-masing menjadi satu baris.
\end{example}

\begin{example}[Modus sebuah hukum Poisson]
\label{ex:b2:genfun:poissonmode}
Di mana $\P(X = k)$ terbesar untuk $X \sim \mathcal
P(\lambda)$? Bobot berurutannya dibandingkan lewat nisbah
\[
\frac{\P(X = k+1)}{\P(X = k)} = \frac{\lambda}{k + 1} ,
\]
yang melampaui $1$ selama $k < \lambda - 1$ lalu jatuh di bawah
$1$ begitu $k > \lambda - 1$: bobotnya naik lalu turun, dengan
modus $\floor\lambda$ (dan seri antara $\lambda - 1$ dan
$\lambda$ ketika $\lambda$ sebuah bilangan bulat: untuk $\lambda = 3$,
$\P(X = 2) = \P(X = 3) = \frac92\eu^{-3} \approx 0.224$).
Uji nisbah atas koefisiennya kerap menjadi jalan yang tercepat
menuju fakta kualitatif tentang sebuah \omterm{def:b2:randomvar:law}{hukum} diskret --- tanpa perlu \omterm{ex:b2:powerseries:fibonacci}{fungsi
pembangkit}, tetapi koefisiennya \emph{adalah} \omterm{ex:b2:powerseries:fibonacci}{fungsi
pembangkitnya}, yang dibaca suku demi suku.
\end{example}

\section{Jumlah peubah yang saling bebas}

\begin{theorem}[Kemultiplikatifan]\label{thm:b2:genfun:product}
Bila $X$ dan $Y$ \omterm{def:b2:randomvar:law}{peubah acak} bernilai-$\N$ yang \omterm{def:b2:proba:independence}{saling bebas}, maka
\[
G_{X + Y}(t) = G_X(t)\,G_Y(t)
\qquad (\abs t \leq 1),
\]
dan menurut induksi $G_{X_1 + \dots + X_n} = \prod_i G_{X_i}$ bagi
$X_1, \dots, X_n$ yang \omterm{def:b2:proba:independence}{saling bebas}.
\end{theorem}

\begin{proof}
Dua bukti, keduanya mendidik. \emph{Lewat \omterm{def:b2:randomvar:expectation}{nilai harapan}:} $t^X$ dan
$t^Y$ merupakan peubah terbatas yang \omterm{def:b2:proba:independence}{saling bebas}, jadi
(\cref{thm:b2:randomvar:product})
\[
G_{X+Y}(t) = \E\bigl(t^{X+Y}\bigr)
= \E\bigl(t^X t^Y\bigr)
= \E\bigl(t^X\bigr)\E\bigl(t^Y\bigr) .
\]
\emph{Lewat hasil kali Cauchy:} \omterm{def:b2:randomvar:law}{hukum} $X + Y$ adalah konvolusi
$\P(X + Y = n) = \sum_{k=0}^n \P(X = k)\P(Y = n - k)$, dan
teorema hasil kali Cauchy bagi deret yang konvergen \omterm{def:b2:series:def}{mutlak}
(\cref{ch:b2:series}) mengalikan kedua deret pangkatnya persis
sepanjang konvolusi itu.
\end{proof}

\begin{example}[Kestabilan hukum klasiknya]\label{ex:b2:genfun:stability}
Binomial \omterm{def:b2:proba:independence}{saling bebas} dengan $p$ yang sama itu menjumlah: $(1 - p + pt)^m(1 -
p + pt)^n = (1 - p + pt)^{m+n}$, jadi $\mathcal{B}(m, p) +
\mathcal{B}(n, p) = \mathcal{B}(m + n, p)$ --- khususnya jumlah
$n$ peubah Bernoulli yang \omterm{def:b2:proba:independence}{saling bebas} itu binomial, yang membuktikan ulang
\omterm{def:b2:randomvar:law}{hukum} banyaknya keberhasilan. Poisson yang \omterm{def:b2:proba:independence}{saling bebas} pun menjumlah:
$e^{\lambda(t-1)}e^{\mu(t-1)} = e^{(\lambda + \mu)(t-1)}$, jadi
$\mathcal{P}(\lambda) + \mathcal{P}(\mu) = \mathcal{P}(\lambda +
\mu)$ --- yakni perhitungan konvolusi pada
\cref{exo:b2:randomvar:2}, kini tanpa perhitungan.
\end{example}

\begin{example}[Dua dadu, satu suku banyak
dikuadratkan]\label{ex:b2:genfun:twodice}
Untuk satu dadu setimbang, $G(t) = \frac{t + t^2 + \dots + t^6}{6}$;
adapun untuk jumlah dua dadu,
\[
G(t)^2 = \frac{1}{36}\bigl(t^2 + 2t^3 + 3t^4 + 4t^5 + 5t^6 +
6t^7 + 5t^8 + 4t^9 + 3t^{10} + 2t^{11} + t^{12}\bigr) :
\]
yakni \omterm{def:b2:randomvar:law}{hukum} segitiga jumlah dadu ($7$ adalah modusnya, dengan
peluang $\frac6{36} = \frac16$), yang terbaca dari sebuah suku banyak
kuadrat yang orang kalikan sekali seumur hidup. Rumus
konvolusinya akan menuntut sebelas hujah pencacahan terpisah;
adapun \omterm{ex:b2:powerseries:fibonacci}{fungsi pembangkitnya} mengerjakan semuanya serentak,
sebab mengalikan suku banyak \emph{adalah} mengonvolusikan koefisien.
Pengalihbahasaan mekanis ini --- \omterm{def:b2:randomvar:law}{hukum} menjadi koefisien, jumlah menjadi
hasil kali --- merupakan seluruh model usaha bab ini, dan
\cref{exo:b2:genfun:11} mendorongnya sampai ke dadu Sicherman
yang mengejutkan.
\end{example}

\begin{example}[Tiga dadu dan sebuah pemerasan
koefisien]\label{ex:b2:genfun:threedice}
Untuk jumlah $S$ tiga dadu setimbang, $\P(S = 10)$ adalah
koefisien $t^{10}$ pada $\bigl(\frac{t + \dots +
t^6}6\bigr)^3$. Faktorkan lalu uraikan dengan deret binomial dan
deret geometrik:
\[
\Bigl(\frac{t(1 - t^6)}{6(1 - t)}\Bigr)^{\!3}
= \frac{t^3}{216}\,\bigl(1 - 3t^6 + 3t^{12} -
t^{18}\bigr)\sum_{j\geq0}\binom{j+2}2t^j .
\]
Koefisien $t^{10}$ menuntut $t^7$ dari hasil kalinya:
$j = 7$ dengan suku $1$, dan $j = 1$ dengan suku
$-3t^6$:
\[
\P(S = 10) = \frac{1}{216}\Bigl(\binom92 -
3\binom32\Bigr) = \frac{36 - 9}{216} = \frac{27}{216} =
\frac18 .
\]
Pencacahan langsung atas $27$ tripelnya rawan keliru; adapun
aljabarnya mekanis dan berskala ke berapa pun banyaknya dadu ---
inklusi--eksklusi yang terlihat pada $(1 - t^6)^3$ sedang
mengerjakan telaah kasusnya secara otomatis.
\end{example}

\begin{example}[Membaca sebuah hukum dari fungsi
pembangkitnya]\label{ex:b2:genfun:identify}
\omterm{def:b2:randomvar:law}{Hukum} mana yang punya $G(t) = \dfrac1{2 - t}$? Uraikan menjadi deret
pangkat:
\[
\frac{1}{2 - t} = \frac12\cdot\frac1{1 - t/2}
= \sum_{k\geq0}\frac{t^k}{2^{k+1}} :
\]
yakni koefisien taknegatif yang berjumlah $G(1) = 1$, jadi ini sebuah
\omterm{def:b2:randomvar:law}{hukum} sejati, $\P(X = k) = 2^{-(k+1)}$ pada $\N$ --- sebuah
\omterm{def:b2:randomvar:law}{hukum} geometrik yang bermula di $0$. Menurut ketunggalannya
(\cref{prop:b2:genfun:radius}), tak ada \omterm{def:b2:randomvar:law}{hukum} lain yang berbagi
$G$ ini. Mengenali \omterm{def:b2:randomvar:law}{hukum} dari \omterm{ex:b2:powerseries:fibonacci}{fungsi pembangkitnya} merupakan
keterampilan yang layak dilatih: begitulah iterat percabangan
kritis $G_n(t) = \frac{n - (n-1)t}{n+1 - nt}$ pada
soal akhir pekan itu tersingkap sebagai \omterm{def:b2:randomvar:law}{hukum} geometrik yang
disyaratkan pada kesintasannya.
\end{example}

\begin{remark}
Kestabilannya berlaku satu arah saja: jumlah Poisson yang \omterm{def:b2:proba:independence}{saling bebas} itu
Poisson, tetapi selisihnya tidak --- $X - Y$ mengambil nilai
negatif, jadi ia sama sekali tak punya \omterm{ex:b2:powerseries:fibonacci}{fungsi pembangkit}, dan \omterm{def:b2:randomvar:law}{hukumnya}
(\omterm{def:b2:randomvar:law}{distribusi} Skellam) terletak di luar perkakas bab ini.
Begitu pula $\mathcal B(m, p) + \mathcal B(n, p')$ dengan
$p \neq p'$ \emph{bukanlah} binomial: hasil kali $(1 - p +
pt)^m(1 - p' + p't)^n$ punya dua letak akar yang berbeda, sedangkan
setiap \omterm{ex:b2:powerseries:fibonacci}{fungsi pembangkit} binomial punya satu akar berulang tunggal. Membaca
kestabilan dari pola akarnya merupakan cuplikan kecil tentang betapa banyak
struktur yang disandikan suku banyaknya.
\end{remark}

\begin{remark}[Tapis akar satuan]
Mengevaluasinya di $-1$ memisahkan yang genap dari yang ganjil; mengevaluasinya di semua
akar satuan ke-$m$ memisahkan setiap kelas sisanya: dengan
$\omega = \eu^{2\iu\pi/m}$,
\[
\P(X \equiv r \bmod m)
= \frac1m\sum_{j=0}^{m-1}\omega^{-jr}\,G_X(\omega^j),
\]
sebab merata-ratakan $\omega^{j(k-r)}$ atas $j$ menghasilkan $1$ bila $k
\equiv r$ dan $0$ bila tidak. Contoh hasilnya: untuk jumlah $S$
dua dadu setimbang, setiap $G(\omega^j) = \frac16\sum_{k=1}^6
\omega^{jk} = -\frac16$ untuk $j \neq 0$ (ketujuh akar
satuan ketujuh berjumlah nol), jadi
\[
\P(7 \mid S) = \frac17\Bigl(1 +
6\cdot\frac1{36}\Bigr) = \frac16 ,
\]
yang membenarkan cacahan pada \cref{ex:b2:genfun:twodice} ---
dan metodenya berskala ke pertanyaan yang tak terjangkau
pencacahan langsung.
\end{remark}

\begin{theorem}[Jumlah acak: kesamaan Wald bagi fungsi
pembangkit]\label{thm:b2:genfun:compound}
Misalkan $(X_k)_{k\geq1}$ peubah bernilai-$\N$ yang \omterm{def:b2:proba:independence}{saling bebas} dengan
\omterm{def:b2:randomvar:law}{hukum} dan \omterm{ex:b2:powerseries:fibonacci}{fungsi pembangkit} $G_X$ yang sama, dan misalkan $N$ sebuah
peubah bernilai-$\N$ yang bebas dari $X_k$, dengan \omterm{ex:b2:powerseries:fibonacci}{fungsi
pembangkit} $G_N$. Maka jumlah acak $S = X_1 + \dots + X_N$ (dengan
$S = 0$ ketika $N = 0$) punya \omterm{ex:b2:powerseries:fibonacci}{fungsi pembangkit}
\[
G_S = G_N \circ G_X .
\]
Khususnya, bila $N$ dan $X_1$ punya \omterm{def:b2:randomvar:expectation}{nilai harapan}, maka $\E(S) =
\E(N)\,\E(X_1)$.
\end{theorem}

\begin{proof}
Syaratkanlah pada $N$ (peluang total,
\cref{thm:b2:proba:bayes}): untuk $\abs t \leq 1$,
\[
G_S(t) = \sum_{n=0}^\infty \P(N = n)\,
\E\bigl(t^{X_1 + \dots + X_n}\bigr)
= \sum_{n=0}^\infty \P(N = n)\,G_X(t)^n
= G_N\bigl(G_X(t)\bigr),
\]
dengan memakai kemultiplikatifannya bagi setiap $n$ yang tetap dan keterjumlahan
seluruh keluarga rangkapnya ($\abs{G_X(t)} \leq 1$). Adapun pertukaran
penjumlahannya adalah Fubini bagi keluarga yang \omterm{def:b2:series:summable}{terjumlahkan}
(\cref{ch:b2:series}). Menurunkannya di $1^-$ dengan aturan rantai
dan \cref{thm:b2:genfun:moments}: $\E(S) = G_N'(G_X(1))\,G_X'(1)
= G_N'(1)G_X'(1) = \E(N)\E(X_1)$.
\end{proof}

\begin{example}[Poisson majemuk: kerugian asuransi
setahun]\label{ex:b2:genfun:compoundpoisson}
Sebuah penanggung menerima $N \sim \mathcal P(\lambda)$ klaim dalam
setahun, dengan setiap klaim berbiaya $X_k$ (satuan bulat, berhukum sama dan
\omterm{def:b2:proba:independence}{saling bebas}, \omterm{ex:b2:powerseries:fibonacci}{fungsi pembangkit} $G_X$, rerata $\mu$, bebas dari $N$). Menurut
\cref{thm:b2:genfun:compound}, kerugian total $S$ punya
\[
G_S(t) = \eu^{\lambda(G_X(t) - 1)},
\qquad
\E(S) = \lambda\mu ,
\]
dan dengan menurunkannya dua kali di $1^-$:
\[
V(S) = \lambda\,G_X''(1) + \lambda^2\mu^2 + \lambda\mu -
(\lambda\mu)^2 = \lambda\,\E(X^2) .
\]
\omterm{def:b2:randomvar:variance}{Ragamnya} melibatkan momen \emph{kedua} sebuah klaim tunggal,
bukan \omterm{def:b2:randomvar:variance}{ragamnya}: jumlah Poisson majemuk merasakan
klaim besar yang sesekali itu dua kali --- sekali lewat berapa banyaknya, sekali
lewat berapa besarnya. Untuk $\lambda = 10$ klaim berhukum geometrik
dengan rerata $2$ ($\E X^2 = 6$): $\E S = 20$, $V(S) = 60$, dan
Chebyshev (\cref{ch:b2:randomvar}) sudah memberikan
marjin solvabilitas yang terpakai. Pola ``jumlah yang terhenti secara acak'' ini
sama dengan pola yang kelak menggerakkan rekursi percabangan pada
\cref{prop:b2:genfun:branching}: komposisi \omterm{ex:b2:powerseries:fibonacci}{fungsi
pembangkit} adalah aljabar populasi acak.
\end{example}

\begin{remark}
Kebebasan $N$ dari sukunya bukanlah hiasan.
Ambil $X_k \in \{0, 2\}$ dengan peluang sama lalu tetapkan $N =
X_1$ (yang terang-terangan bergantung): maka $S = X_1 + \dots + X_N$ bernilai
$0$ ketika $X_1 = 0$, dan $2 + X_2$ ketika $X_1 = 2$, jadi $\E(S) =
\frac12(2 + 1) = \frac32$, sedangkan $\E(N)\E(X_1) = 1\cdot1 =
1$: kesamaan Wald gagal. Ketika banyaknya suku
dibiarkan \emph{bereaksi} terhadap sukunya sendiri, struktur
hasil kali yang rapi itu runtuh --- adapun teori \omterm{def:b2:metric:complete}{lengkap} kaidah
``penghentian'' semacam itu merupakan bab martingal pada jilid
Tahun ke-3.
\end{remark}

\section{Hampiran Poisson}

\begin{theorem}[Hukum kejadian langka]\label{thm:b2:genfun:rareevents}
Misalkan $X_n \sim \mathcal{B}(n, p_n)$ dengan $n\,p_n \to \lambda > 0$.
Maka untuk setiap $k \in \N$:
\[
\P(X_n = k)
\xrightarrow[n\to\infty]{}
e^{-\lambda}\frac{\lambda^k}{k!} :
\]
yakni \omterm{def:b2:randomvar:law}{hukum} binomial bagi banyak \omterm{def:b2:proba:space}{kejadian} langka yang \omterm{def:b2:proba:independence}{saling bebas} konvergen ke
\omterm{def:b2:randomvar:law}{hukum} Poisson berparameter $\lambda$.
\end{theorem}

\begin{proof}
Perhitungan langsung dengan $p_n = \frac{\lambda_n}{n}$, $\lambda_n
\to \lambda$:
\[
\P(X_n = k)
= \binom nk p_n^k(1 - p_n)^{n-k}
= \frac{n(n-1)\cdots(n-k+1)}{n^k}\cdot
\frac{\lambda_n^k}{k!}\,
\bigl(1 - \tfrac{\lambda_n}{n}\bigr)^{n-k} .
\]
Ketika $n \to \infty$ dengan $k$ tetap: faktor pertamanya menuju $1$
(hasil kali $k$ faktor $\to 1$); $\lambda_n^k \to \lambda^k$;
dan $\bigl(1 - \frac{\lambda_n}{n}\bigr)^{n-k} =
\exp\bigl((n-k)\ln(1 - \frac{\lambda_n}{n})\bigr) \to
e^{-\lambda}$ sebab $(n - k)\ln\bigl(1 - \frac{\lambda_n}{n}\bigr)
\sim -\lambda_n \to -\lambda$
(\cref{ch:b2:comparison}). Sebagai gantinya, pada tataran
\omterm{ex:b2:powerseries:fibonacci}{fungsi pembangkit}: $G_{X_n}(t) = \bigl(1 +
\frac{\lambda_n(t-1)}{n}\bigr)^n \to e^{\lambda(t - 1)} =
G_{\mathcal{P}(\lambda)}(t)$ untuk setiap $t \in [0, 1]$ yang tetap ---
yakni kekonvergenan \omterm{ex:b2:powerseries:fibonacci}{fungsi pembangkit}, yang (bagi peubah
bernilai-$\N$) setara dengan kekonvergenan setiap $\P(X_n = k)$;
lihat \cref{exo:b2:genfun:9}.
\end{proof}

\begin{remark}
Itulah sebabnya \omterm{def:b2:randomvar:law}{hukum} Poisson dipakai untuk memodelkan \omterm{def:b2:proba:space}{kejadian} langka --- salah ketik per
halaman, peluruhan radioaktif per detik, kecelakaan per hari di sebuah
persimpangan: setiap kesempatannya nyaris dapat diabaikan, kesempatannya
banyak, dan hanya laju reratanya $\lambda$ yang bertahan pada limitnya.
\end{remark}

\begin{example}[Menyaksikan limit Poisson
konvergen]\label{ex:b2:genfun:poissonnumerics}
Tetapkan $\lambda = 2$ lalu ambil $X_n \sim \mathcal B(n, 2/n)$. Adapun
peluang tanpa \omterm{def:b2:proba:space}{kejadian} persis $\P(X_n = 0) = (1 - 2/n)^n$:
\[
n = 10:\ 0.107, \qquad
n = 20:\ 0.122, \qquad
n = 50:\ 0.130, \qquad
n = 100:\ 0.133,
\]
terhadap limitnya $\eu^{-2} \approx 0.135$. Kekonvergenannya
monoton dan berlaju $O(1/n)$ --- dengan menguraikannya, $(1 -
2/n)^n = \eu^{-2}\bigl(1 - \tfrac2n + O(n^{-2})\bigr)$ ---
jadi untuk $n$ yang bernilai ratusan model Poissonnya sudah
tepat sampai digit ketiga. Itulah isi praktis
\omterm{def:b2:randomvar:law}{hukum} \omterm{def:b2:proba:space}{kejadian} langka: pemodelnya tak pernah mengetahui $n$ dan
$p$ secara terpisah (berapa banyak kesempatan mikro salah ketik yang
dimuat satu halaman?), melainkan hanya hasil kalinya $\lambda$, dan
\omterm{def:b2:randomvar:law}{hukum} limitnya dengan murah hati tak bergantung pada apa pun yang lain.
\end{example}

\section{Proses percabangan}

Tinjaulah sebuah populasi yang bermula dari satu leluhur; setiap individu,
secara \omterm{def:b2:proba:independence}{saling bebas}, punya banyak anak yang acak dengan \omterm{def:b2:randomvar:law}{hukum} $(p_k)_{k
\in \N}$ dan \omterm{ex:b2:powerseries:fibonacci}{fungsi pembangkit} $G$ (yakni \emph{\omterm{def:b2:randomvar:law}{distribusi}
keturunan}). Misalkan $Z_n$ ukuran generasi $n$ ($Z_0 =
1$), dan misalkan $m = G'(1) = \E(Z_1)$ rerata banyaknya keturunan.

\begin{proposition}\label{prop:b2:genfun:branching}
\omterm{ex:b2:powerseries:fibonacci}{Fungsi pembangkit} $Z_n$ adalah iterat ke-$n$
$G_{Z_n} = G \circ G \circ \dots \circ G$ (sebanyak $n$ kali), dan
peluang kepunahan $q_n = \P(Z_n = 0)$ memenuhi
\[
q_0 = 0, \qquad q_{n+1} = G(q_n),
\]
lalu naik menuju peluang $q$ kepunahan pada akhirnya, yang
merupakan titik tetap $G$.
\end{proposition}

\begin{proof}
Generasi $n + 1$ merupakan jumlah acak keturunan dari $Z_n$
anggota generasi $n$, dengan cacahnya \omterm{def:b2:proba:independence}{saling bebas} satu sama lain
dan bebas dari $Z_n$: \cref{thm:b2:genfun:compound} memberi $G_{Z_{n+1}} =
G_{Z_n} \circ G$, dan induksi dari $G_{Z_0}(t) = t$ menghasilkan
iterat $n$ lapisnya --- yang, menurut keasosiatifan komposisinya, dapat
pula dibaca sebagai $G_{Z_{n+1}} = G \circ G_{Z_n}$. Dengan mengevaluasi
bentuk kedua ini di $0$: $q_{n+1} = G_{Z_{n+1}}(0) =
G\bigl(G_{Z_n}(0)\bigr) = G(q_n)$. \omterm{def:b2:proba:space}{Kejadian} $\{Z_n = 0\}$ naik (populasi yang punah tetap
punah), jadi $q_n \uparrow q = \P\bigl(\bigcup_n\{Z_n = 0\}\bigr)$
menurut kekontinuan monoton (\cref{thm:b2:proba:continuity}), dan
kekontinuan $G$ pada $[0, 1]$ mengubah $q_{n+1} = G(q_n)$ menjadi $q =
G(q)$ pada limitnya.
\end{proof}

\begin{example}[Menyaksikan kepunahan konvergen]
\label{ex:b2:genfun:cobwebnumerics}
Untuk \omterm{def:b2:randomvar:law}{hukum} keturunan $(p_0, p_1, p_2) = (\tfrac14,
\tfrac14, \tfrac12)$ pada \cref{ex:b2:genfun:branchingexample},
$G(t) = \tfrac14 + \tfrac14t + \tfrac12t^2$ dan iterasi
$q_{n+1} = G(q_n)$ memberi
\[
q_1 = 0.25, \quad q_2 = 0.34375, \quad q_3 \approx 0.39502,
\quad q_4 \approx 0.42678, \quad q_5 \approx 0.44776,
\]
yang memanjat menuju peluang kepunahan $q = \tfrac12$.
Selisihnya $q - q_n$ adalah $0.25$, $0.156$, $0.105$, $0.073$,
$0.052$: masing-masing kira-kira $\tfrac34$ kali sebelumnya, dan
memang teorema nilai rata-rata memberi $q - q_{n+1} = G'(c_n)(q
- q_n)$ dengan $G'(q) = \tfrac14 + q = \tfrac34$. Dua pelajarannya: sebuah
garis keluarga yang masih hidup pada generasi $n$ punya, terpasang dalam
perhitungan yang sama, peluang $q - q_n$ untuk punah
kemudian; dan laju kekonvergenan tangga pada
gambar di bawah adalah turunannya di titik tetapnya --- adapun
soal akhir pekannya mengubah kedua pengamatan itu menjadi teorema.
\end{example}

\begin{theorem}[Kriteria kepunahan]\label{thm:b2:genfun:extinction}
Andaikan $p_1 \neq 1$. Peluang kepunahan $q$ merupakan titik tetap
\emph{terkecil} $G$ pada $\intcc{0}{1}$, dan:
\begin{itemize}
\item bila $m \leq 1$ (subkritis atau kritis), maka $q = 1$: kepunahannya
pasti;
\item bila $m > 1$ (superkritis), maka $q < 1$: populasinya sintas
selamanya dengan peluang positif $1 - q$.
\end{itemize}
\end{theorem}

\begin{proof}
$G$ \omterm{def:b2:affine:convex}{cembung} pada $\intcc{0}{1}$ (deret pangkat berkoefisien
taknegatif: $G'' \geq 0$), tidak turun, dengan $G(1) = 1$.

\emph{Titik tetap terkecil:} ambil $r \in \intcc{0}{1}$ sebarang titik
tetap. Maka $q_0 = 0 \leq r$, dan secara induktif $q_{n+1} = G(q_n)
\leq G(r) = r$ (kemonotonannya): jadi $q = \lim q_n \leq r$.

\emph{Kasus $m \leq 1$:} andaikan $r < 1$ sebuah titik tetap. Menurut
teorema nilai rata-rata pada $[r, 1]$, ada $c \in \intoo{r}{1}$ dengan
$G'(c) = \frac{G(1) - G(r)}{1 - r} = \frac{1 - r}{1 - r} = 1$. Padahal
$G'$ tidak turun (kecembungannya) dengan $\lim_{t\to1^-}G'(t) = m
\leq 1$, jadi $G' \leq 1$ pada $\intoo{0}{1}$; kesamaan $G'(c) =
1$ lalu memaksa $G'$ konstan bernilai $1$ pada $\intco{c}{1}$,
sehingga $G'' = \sum n(n-1)p_nt^{n-2} \equiv 0$ di sana. Deret pangkat
berkoefisien taknegatif yang lenyap pada sebuah selang punya semua
koefisien itu nol: $p_n = 0$ untuk $n \geq 2$, jadi $G(t) = p_0
+ p_1t$ dan $1 = G'(c) = p_1$ --- yang bertentangan dengan hipotesis
$p_1 \neq 1$. Jadi $1$ satu-satunya titik tetapnya: $q = 1$.

\emph{Kasus $m > 1$:} di dekat $1$, $G(t) - t$ berturunan $G'(t) -
1 \to m - 1 > 0$ ketika $t \to 1^-$, jadi $G(t) - t < G(1) - 1 = 0$ pada
suatu selang $\intoo{1 - \delta}{1}$: fungsi \omterm{def:b2:metric:continuity}{kontinu}
$G(t) - t$ bernilai $\geq 0$ di $t = 0$ ($G(0) = p_0 \geq 0$) dan $< 0$
tepat di bawah $1$, jadi ia lenyap di suatu $r < 1$ (teorema nilai
antara). Titik tetap terkecilnya lalu $q \leq r < 1$.
\end{proof}

\begin{figure}[ht]
\centering
% Subcritical: G(t) = 0.5 + 0.1 t + 0.4 t^2 => m = 0.9 <= 1
\begin{tikzpicture}[scale=2.9]
  \draw[->] (0,0) -- (1.15,0) node[below] {$t$};
  \draw[->] (0,0) -- (0,1.15);
  \draw[dashed] (0,0) -- (1,1);
  \draw[thick, ocBlue, domain=0:1, smooth, samples=40]
    plot (\x, {0.5 + 0.1*\x + 0.4*\x*\x});
  \fill (1,1) circle (0.02);
  \node[below, font=\small] at (0.55,-0.06)
    {$m \leq 1$: $q = 1$};
  % cobweb from 0
  \draw[ocRed, ->] (0,0) -- (0,0.5);
  \draw[ocRed, ->] (0,0.5) -- (0.5,0.5);
  \draw[ocRed, ->] (0.5,0.5) -- (0.5,0.65);
  \draw[ocRed, ->] (0.5,0.65) -- (0.65,0.65);
  \draw[ocRed, ->] (0.65,0.65) -- (0.65,0.734);
  \draw[ocRed, ->] (0.65,0.734) -- (0.734,0.734);
  \node[font=\small, ocBlue] at (0.32,0.75) {$G$};
\end{tikzpicture}
\qquad
% Supercritical: G(t) = 0.25 + 0.15 t + 0.6 t^2 => m = 1.35 > 1
\begin{tikzpicture}[scale=2.9]
  \draw[->] (0,0) -- (1.15,0) node[below] {$t$};
  \draw[->] (0,0) -- (0,1.15);
  \draw[dashed] (0,0) -- (1,1);
  \draw[thick, ocBlue, domain=0:1, smooth, samples=40]
    plot (\x, {0.25 + 0.15*\x + 0.6*\x*\x});
  \fill (1,1) circle (0.02);
  % fixed point q: 0.6q^2 - 0.85q + 0.25 = 0 -> q = 5/12 ~ 0.4167
  \fill[ocRed] (0.4167,0.4167) circle (0.02)
    node[below right, font=\small] {$q$};
  \node[below, font=\small] at (0.55,-0.06)
    {$m > 1$: $q < 1$};
  % cobweb from 0
  \draw[ocRed, ->] (0,0) -- (0,0.25);
  \draw[ocRed, ->] (0,0.25) -- (0.25,0.25);
  \draw[ocRed, ->] (0.25,0.25) -- (0.25,0.325);
  \draw[ocRed, ->] (0.25,0.325) -- (0.325,0.325);
  \draw[ocRed, ->] (0.325,0.325) -- (0.325,0.3623);
  \draw[ocRed, ->] (0.325,0.3623) -- (0.3623,0.3623);
  \node[font=\small, ocBlue] at (0.28,0.72) {$G$};
\end{tikzpicture}
\caption{Peluang kepunahan sebagai iterasi titik tetap
$q_{n+1} = G(q_n)$ yang bermula di $q_0 = 0$ (tangga merah). Kiri: sebuah
\omterm{def:b2:randomvar:law}{hukum} keturunan subkritis --- kurva \omterm{def:b2:affine:convex}{cembungnya} tetap di atas
diagonalnya, dan iterasinya memanjat ke titik tetap tunggal $1$.
Kanan: sebuah \omterm{def:b2:randomvar:law}{hukum} superkritis --- kurvanya memotong diagonalnya di
$q < 1$, tempat iterasinya berhenti: kesintasannya berpeluang $1 -
q > 0$.}
\label{fig:b2:genfun:branching}
\end{figure}

\begin{remark}[Cara membaca sarang laba-labanya]
Pada gambarnya, langkah tegak menerapkan $G$ (dari $(q_n, q_n)$
naik ke $(q_n, G(q_n))$), sedangkan langkah mendatar menuju diagonalnya
mengubah keluaran menjadi masukan: tangganya \emph{adalah}
rekursi $q_{n+1} = G(q_n)$. Kecembungan $G$ dan $G(1) = 1$
menyisakan dua geometri saja. Entah kurvanya tetap di atas
diagonalnya pada $\intco01$ (rerata $m \leq 1$): tangganya tak punya
tempat berhenti sebelum $1$. Entah kurvanya memotong di suatu $q <
1$ ($m > 1$): tangganya terperangkap di bawah perpotongannya
lalu konvergen ke sana, dengan laju geometrik $G'(q) < 1$
yang dikuantifikasi pada \cref{ex:b2:genfun:cobwebnumerics}. Seluruh
telaah teorema kepunahannya terlihat pada satu
gambar ini --- itulah sebabnya ia layak digambar sebelum
dihitung.
\end{remark}

\begin{example}\label{ex:b2:genfun:branchingexample}
\omterm{def:b2:randomvar:law}{Hukum} keturunan: tanpa anak, satu anak, dua anak dengan
peluang $\frac14, \frac14, \frac12$. Maka $m = \frac14 + 1 =
\frac54 > 1$ dan $G(t) = \frac14 + \frac14 t + \frac12 t^2$. Titik
tetapnya: $\frac12 t^2 - \frac34 t + \frac14 = 0$, yakni $2t^2 - 3t
+ 1 = (2t - 1)(t - 1) = 0$: $q = \frac12$. Garis keluarganya punah
dengan peluang $\frac12$ --- dan dengan peluang $\frac12$
ia hidup selamanya.
\end{example}

\begin{remark}[Perspektif di dalam jilid ini]
Bab ini merupakan persimpangan buku ini, dan setiap bahannya
datang dari tempat yang bernama: aljabar deretnya dari
\cref{ch:b2:series} dan \cref{ch:b2:powerseries}, peluangnya
dari \cref{ch:b2:proba} (kekontinuan monoton membuktikan
$q_n \uparrow q$) dan \cref{ch:b2:randomvar} ($G_X =
\E(t^X)$ sebuah \omterm{def:b2:randomvar:expectation}{nilai harapan}, kemultiplikatifannya adalah teorema
hasil kali), kecembungannya dari \cref{ch:b2:realfun} sampai
\cref{ch:b2:affine}. Bahkan kepatologian ekor beratnya pun \omterm{def:b2:metric:connected}{terhubung}:
peubah St Petersburg pada bab sebelumnya punya
$G(t) = \sum_k2^{-k}t^{2^k}$, yakni deret yang konvergen sempurna
pada $\intcc01$ yang turunannya di $1^-$ divergen ---
rerata tak hingga, terlihat sekilas pandang. Satu objek, segala perkakas
tahun ini: sebuah bab penutup yang pantas.
\end{remark}

\begin{remark}[Jebakan yang lazim]
(i) \omterm{ex:b2:powerseries:fibonacci}{Fungsi pembangkit} hanya berlaku bagi peubah bernilai-$\N$:
bagi peubah bertanda atau tak bulat, objek $\E(t^X)$
kehilangan struktur deret pangkatnya (Tahun ke-3 menggantinya dengan
transformasi yang disesuaikan dengan $\R$). (ii) Pemeriksaan kewarasan pertama atas
setiap $G$ yang terhitung adalah $G(1) = 1$; yang kedua adalah bahwa
koefisiennya taknegatif --- koefisien negatif berarti
kekeliruan aljabar, bukan \omterm{def:b2:randomvar:law}{hukum} baru. (iii) Pada jumlah acak,
urutan komposisinya penting: $G_S = G_N \circ G_X$, dengan
fungsi \emph{luarnya} mencacah sukunya; menyusunnya
dengan cara sebaliknya tak bermakna ($G_X \circ G_N$ akan mencacah benda
dari benda). (iv) Kemultiplikatifannya perlu kesalingbebasan dan
sumber keacakan yang berbeda: $G_{2X}(t) = G_X(t^2)$, bukan
$G_X(t)^2$. (v) Menurunkannya di $1$ merupakan operasi perbatasan:
ketika jari-jarinya persis $1$, seperti pada
\cref{ex:b2:genfun:heavytail}, $G'(1^-)$ boleh jadi tak hingga,
dan perumusan limit monoton pada teorema momennya bukanlah
kerewelan bertele-tele melainkan pernyataan yang jujur.
\end{remark}

\section*{Menutup jilid ini}

\omterm{ex:b2:powerseries:fibonacci}{Fungsi pembangkit} merupakan objek penutup yang pantas bagi buku ini:
ia sekaligus sebuah deret pangkat (\cref{ch:b2:powerseries}), sebuah
perkakas keluarga yang \omterm{def:b2:series:summable}{terjumlahkan} (\cref{ch:b2:series}), sebuah \omterm{def:b2:randomvar:expectation}{nilai harapan}
(\cref{ch:b2:randomvar}), sebuah fungsi \omterm{def:b2:affine:convex}{cembung} yang geometrinya memutuskan
kepunahan (\cref{ch:b2:realfun}), dan sebuah iterasi titik tetap
(\cref{ch:b2:metric}). Matematika Tahun ke-2 adalah satu pokok bahasan.
Adapun jilid Tahun ke-3 akan membuka pintu yang sengaja dibiarkan tertutup
di sini: integral Lebesgue (yang melunasi teorema kekonvergenan
terdominasi pada \cref{ch:b2:integration}), peluang berteori ukuran
pada ruang yang takterbilang, dan bukti \omterm{def:b2:metric:complete}{lengkap} teorema fungsi
invers (\cref{ch:b2:diffcalc}) dalam latar geometri \omterm{def:b2:diffcalc:differential}{diferensial}.

\section{Latihan}

\begin{exercise}[$\star$]\label{exo:b2:genfun:1}
Hitunglah \omterm{ex:b2:powerseries:fibonacci}{fungsi pembangkit} hukum seragam pada $\{1, 2,
\dots, 6\}$ (sebuah dadu setimbang). Tunjukkan bahwa jumlah dua dadu setimbang
\emph{tak mungkin} seragam pada $\{2, \dots, 12\}$: faktorkan $G_{X+Y}$
lalu cacah akarnya. \emph{(Jumlah yang seragam akan memaksa $G_X(t)G_Y(t) =
\frac{t^2}{11}\sum_{k=0}^{10}t^k$, yang akar taknolnya adalah
akar satuan ke-$11$ selain $1$ --- dan tak satu pun real ---
sedangkan $G_X/t$ dan $G_Y/t$ merupakan suku banyak real berderajat $5$, yang
masing-masing memiliki sedikitnya satu akar real.)}
\end{exercise}

\begin{exercise}[$\star$]\label{exo:b2:genfun:2}
Dengan memakai \omterm{ex:b2:powerseries:fibonacci}{fungsi pembangkit}, pulihkan $\E$ dan $V$ bagi \omterm{def:b2:randomvar:law}{hukum} binomial
dan \omterm{def:b2:randomvar:law}{hukum} geometrik (\cref{thm:b2:genfun:moments}).
\end{exercise}

\begin{exercise}[$\star$]\label{exo:b2:genfun:3}
Dua dadu yang dicurangi: mungkinkah mencurangi dua dadu (secara \omterm{def:b2:proba:independence}{saling bebas},
seidentik atau tidak) sehingga jumlahnya seragam pada $\{2, \dots,
12\}$? \emph{(Halangan pemfaktoran yang sama seperti pada
\cref{exo:b2:genfun:1}: jawabannya tidak, bahkan dengan kecurangan yang berbeda,
sebab setiap faktor $G_X(t)/t$ berderajat ganjil $5$,
sehingga punya akar real, sedangkan sasarannya tak punya.)}
\end{exercise}

\begin{exercise}[$\star\star$]\label{exo:b2:genfun:4}
Misalkan $X_1, X_2, \dots$ Bernoulli $\mathcal{B}(p)$ yang \omterm{def:b2:proba:independence}{saling bebas}
dan $N \sim \mathcal{P}(\lambda)$ bebas darinya. Tunjukkan,
lewat \cref{thm:b2:genfun:compound}, bahwa $S = X_1 + \dots + X_N
\sim \mathcal{P}(\lambda p)$: sebanyak Poisson benda, yang masing-masing disimpan
dengan peluang $p$, menyisakan sebanyak Poisson pula --- yakni \emph{penipisan}.
Hitunglah pula \omterm{def:b2:randomvar:law}{hukum} cacah yang terbuang lalu kagumilah: ia
$\mathcal{P}(\lambda(1-p))$, dan dapat ditunjukkan bahwa ia bebas
dari $S$.
\end{exercise}

\begin{exercise}[$\star\star$]\label{exo:b2:genfun:5}
(Binomial negatif) Misalkan $T_r$ banyaknya lemparan untuk memperoleh
$r$ gambar (peluang gambar $p$). Tuliskan $T_r$ sebagai jumlah $r$
peubah geometrik yang \omterm{def:b2:proba:independence}{saling bebas}, lalu simpulkan
\[
G_{T_r}(t) = \Bigl(\frac{pt}{1 - (1-p)t}\Bigr)^{r},
\qquad
\E(T_r) = \frac rp,
\qquad
V(T_r) = \frac{r(1-p)}{p^2},
\]
lalu uraikan $G_{T_r}$ untuk memperoleh $\P(T_r = n) = \binom{n-1}{r-1}
p^r(1-p)^{n-r}$.
\end{exercise}

\begin{exercise}[$\star\star$]\label{exo:b2:genfun:6}
Untuk \omterm{def:b2:randomvar:law}{hukum} keturunan $p_0 = \frac18$, $p_1 = \frac38$, $p_2 =
\frac38$, $p_3 = \frac18$: hitunglah $m$, putuskanlah
kesuperkritisannya, lalu hitunglah peluang kepunahan $q$
secara persis. \emph{(Faktorkan keluar akar $t = 1$ dari $G(t) - t$.)}
\end{exercise}

\begin{exercise}[$\star\star\star$]\label{exo:b2:genfun:7}
(Total keturunan) Pada sebuah \omterm{pb:b2:genfun:1}{proses percabangan} subkritis ($m < 1$), misalkan
$Y = \sum_{n\geq0} Z_n$ jumlah seluruh individu yang pernah
lahir. Tunjukkan $\E(Y) = \sum_n m^n = \frac{1}{1 - m}$ (sahkanlah
pertukaran penjumlahannya), lalu buktikan bahwa \omterm{ex:b2:powerseries:fibonacci}{fungsi pembangkit}
$H = G_Y$ memenuhi persamaan fungsional $H(t) = t\,G(H(t))$.
\emph{(Leluhurnya, ditambah total keturunan masing-masing
anaknya, yang merupakan salinan bebas dari $Y$.)}
\end{exercise}

\begin{exercise}[$\star\star\star$]\label{exo:b2:genfun:8}
Misalkan $X$ berfungsi pembangkit $G$ dengan \omterm{def:b2:powerseries:radius}{jari-jari kekonvergenan} $>
1$. Buktikanlah \emph{batas ekor eksponensialnya}: ada $C > 0$ dan
$\rho \in \intoo{0}{1}$ dengan $\P(X \geq n) \leq C\rho^n$.
\emph{(Markov yang diterapkan pada $t^X$ untuk sebuah $t > 1$ yang tetap di dalam
cakramnya.)} Sebaliknya, tunjukkan bahwa bila $\P(X \geq n) \leq C\rho^n$ dengan
$\rho < 1$, maka jari-jari $G$ adalah $\geq 1/\rho > 1$.
\end{exercise}

\begin{exercise}[$\star\star\star$]\label{exo:b2:genfun:9}
(Teorema kekontinuan, kasus dasarnya) Misalkan $X, X_1, X_2, \dots$
bernilai-$\N$ dengan $G_{X_n}(t) \to G_X(t)$ untuk setiap $t \in
\intco{0}{1}$. Tunjukkan bahwa $\P(X_n = k) \to \P(X = k)$ untuk setiap
$k$. \emph{(Induksi atas $k$: untuk $k = 0$ ambil $t \to 0$ ---
dengan hati-hati: tetapkan $t$ yang kecil, pakailah $\abs{\P(X_n = 0) - G_{X_n}(t)}
\leq \frac{t}{1-t}$, yang sah sebab ekornya $\sum_{j \geq 1}p_jt^j
\leq \frac{t}{1 - t}$; lalu diagonalkan. Untuk langkah induksinya,
tinjaulah $\frac{G(t) - \P(X = 0)}{t}$, yakni \omterm{ex:b2:powerseries:fibonacci}{fungsi pembangkit}
sebuah \omterm{def:b2:randomvar:law}{hukum} yang tergeser.)}
\end{exercise}

\begin{exercise}[$\star$]\label{exo:b2:genfun:10}
(Kiat keparitasan) Tunjukkan bahwa untuk peubah bernilai-$\N$ bernama $X$,
\[
\P(X \text{ genap}) = \frac{1 + G_X(-1)}{2} ,
\]
lalu hitunglah peluang ini untuk $X \sim \mathcal P(\lambda)$
dan $X \sim \mathcal B(n, p)$. Apakah makna $G_X(-1) \to 0$
secara peluang?
\end{exercise}

\begin{exercise}[$\star\star$]\label{exo:b2:genfun:11}
(Dadu Sicherman) Sahkanlah pemfaktoran \omterm{ex:b2:powerseries:fibonacci}{fungsi pembangkit}
dadu setimbangnya
\[
\frac{t + t^2 + \dots + t^6}{6}
= \frac{t\,(1 + t)(1 + t + t^2)(1 - t + t^2)}{6},
\]
lalu tunjukkan bahwa kedua dadu bermuka $\{1, 2, 2, 3, 3, 4\}$
dan $\{1, 3, 4, 5, 6, 8\}$ punya \omterm{ex:b2:powerseries:fibonacci}{fungsi pembangkit}
$\frac{t(1+t)(1+t+t^2)}6$ dan
$\frac{t(1+t)(1+t+t^2)(1-t+t^2)^2}6$, yang hasil kalinya sama dengan hasil kali
dua dadu baku: dadu eksotis ini menghasilkan setiap total $2,
\dots, 12$ dengan peluang yang persis sama dengan peluang bakunya.
\end{exercise}

\begin{exercise}[$\star\star\star$]\label{exo:b2:genfun:12}
(Menunggu dua gambar berturut-turut) Sekeping koin dengan peluang gambar
$p$ dilempar sampai muncul dua gambar berurutan; misalkan $T$
banyaknya lemparan (permainan pada \cref{exo:b2:proba:6}).
Dengan menyaratkan pada lemparan pertamanya, turunkan sistem linear bagi
\omterm{ex:b2:powerseries:fibonacci}{fungsi pembangkit} dari keadaan ``belum ada gambar''
dan ``sudah satu gambar'', lalu simpulkan
\[
G_T(t) = \frac{p^2t^2}{1 - qt - pqt^2}
\qquad (q = 1 - p);
\]
periksalah $G_T(1) = 1$ dan $\E(T) = \dfrac{1 + p}{p^2}$ ($= 6$
untuk koin yang setimbang).
\end{exercise}

\section{Soal: proses Galton--Watson, terpecahkan}

\begin{problem}[{Soal akhir pekan --- laju tumbuh, pemecahan
persis, total keturunan, dan taksiran kritis
Kolmogorov}]\label{pb:b2:genfun:1}
Kriteria kepunahan (\cref{thm:b2:genfun:extinction})
membelah \omterm{pb:b2:genfun:1}{proses percabangan}\index{proses percabangan} menjadi
subkritis, kritis, dan superkritis --- tetapi ia tak berkata
apa pun tentang \emph{lajunya}: seberapa cepat garis yang tertakdir punah itu
mati, seberapa besar garis yang sintas itu tumbuh. Soal ini menghitungnya. Kita pertahankan
notasi babnya: \omterm{def:b2:randomvar:law}{hukum} keturunan $(p_k)$ dengan \omterm{ex:b2:powerseries:fibonacci}{fungsi pembangkit} $G$,
rerata $m = G'(1)$, ukuran generasi $Z_n$ ($Z_0 = 1$), iterat
$G_n = G_{Z_n}$, peluang kepunahan $q_n = \P(Z_n = 0)
\uparrow q$; kita selalu mengandaikan $p_1 \neq 1$ dan, di tempat momen
kedua muncul, $G''(1) < \infty$, lalu kita tuliskan $\sigma^2 =
V(Z_1)$.

\medskip
\noindent\textbf{Bagian I --- Momen generasinya.}
\begin{enumerate}
  \item Tunjukkan $\E(Z_n) = m^n$ \emph{(aturan rantai pada $G_n = G
        \circ G_{n-1}$ di $1^-$, dengan memakai $G_{n-1}(1) = 1$ dan
        \cref{thm:b2:genfun:moments})}.
  \item Tegakkan rekursi $G_n''(1) =
        G''(1)\,m^{2(n-1)} + m\,G_{n-1}''(1)$ lalu pecahkan:
        $G_n''(1) = G''(1)\,m^{n-1}\dfrac{m^n - 1}{m - 1}$
        untuk $m \neq 1$, dan $G_n''(1) = n\,G''(1)$ untuk $m =
        1$.
  \item Simpulkan
        \[
        V(Z_n) = \sigma^2m^{n-1}\,\frac{m^n - 1}{m - 1}
        \quad (m \neq 1),
        \qquad
        V(Z_n) = n\,\sigma^2 \quad (m = 1).
        \]
  \item (Laju subkritis, batas atas) Untuk $m < 1$, tunjukkan
        $\P(Z_n > 0) \leq m^n$ \emph{(Markov pada $Z_n$ yang
        bernilai bulat)}: kepunahannya pasti dengan
        laju geometrik --- yakni penghalusan kuantitatif atas
        kriteria babnya.
  \item (Laju subkritis, batas bawah) Dengan memakai
        Cauchy--Schwarz pada $Z_n\mathbf 1_{Z_n > 0}$, tunjukkan
        \[
        \P(Z_n > 0) \geq \frac{\E(Z_n)^2}{\E(Z_n^2)}
        \geq c\,m^{n}
        \quad\text{dengan}\quad
        c = \Bigl(\frac{\sigma^2}{m(1-m)} + 1\Bigr)^{-1} :
        \]
        laju geometriknya $m^n$ persis sampai konstantanya.
\end{enumerate}

\medskip
\noindent\textbf{Bagian II --- Keluarga geometrik, terpecahkan
persis.} Ambil \omterm{def:b2:randomvar:law}{hukum} keturunan yang geometrik pada $\N$: $p_k =
qp^k$ ($k \geq 0$), dengan $0 < p < 1$, $q = 1 - p$.
\begin{enumerate}[resume]
  \item Hitunglah $G(t) = \dfrac{q}{1 - pt}$ dan $m = \dfrac
        pq$; lalu letakkan ketiga rezimnya dalam $p$.
  \item Pecahkan $G(t) = t$: tunjukkan bahwa titik tetapnya $1$ dan
        $q/p = 1/m$, lalu pulihkan peluang kepunahannya
        $q_{\mathrm{ext}} = \min(1, 1/m)$.
  \item Buktikan menurut induksi \omterm{def:b2:multint:exact}{bentuk tertutupnya}
        \[
        q_n = \frac{m^n - 1}{m^{n+1} - 1} \quad (m \neq 1),
        \qquad
        q_n = \frac{n}{n+1} \quad (m = 1).
        \]
  \item Simpulkan laju persisnya: $1 - q_n \sim (1 - m)\,m^n$
        pada kasus subkritisnya, dan $q_{\mathrm{ext}} - q_n
        \sim \dfrac{m - 1}{m^{2}}\cdot m^{-n}$ pada kasus
        superkritisnya; periksalah bahwa nisbah pengerutan
        superkritisnya adalah $G'(q_{\mathrm{ext}}) = 1/m$.
  \item Kasus kritisnya ($p = \tfrac12$): hitunglah $\sigma^2 =
        2$ lalu catat $1 - q_n = \frac1{n+1}$: kesintasannya meluruh
        seperti $\frac1n$ --- tak geometrik dan tak \omterm{def:b2:series:summable}{terjumlahkan}.
  \item Masih kritis: buktikan menurut induksi iterat penuhnya
        \[
        G_n(t) = \frac{n - (n-1)t}{n + 1 - nt},
        \]
        lalu simpulkan bahwa dengan disyaratkan pada kesintasannya, $Z_n$
        geometrik pada $\N^*$ dengan parameter $\frac1{n+1}$:
        \[
        \P(Z_n = k \mid Z_n > 0) = \frac1{n+1}
        \Bigl(\frac{n}{n+1}\Bigr)^{k-1},
        \qquad
        \E(Z_n \mid Z_n > 0) = n + 1 .
        \]
        Garis rata-ratanya mati, tetapi garis yang sintas berukuran
        orde $n$.
\end{enumerate}

\medskip
\noindent\textbf{Bagian III --- Total keturunan.} Misalkan $Y =
\sum_{n\geq0}Z_n \in \N^* \cup \{\infty\}$ jumlah seluruh
individu yang pernah lahir, dan $H(t) =
\sum_{k\geq1}\P(Y = k)t^k$.
\begin{enumerate}[resume]
  \item Sahkanlah $\P(Y < \infty) = q_{\mathrm{ext}}$, lalu
        ingat kembali dari \cref{exo:b2:genfun:7} persamaan
        fungsional $H(t) = t\,G(H(t))$ (yang penurunannya tak
        memakai $m < 1$).
  \item (Percabangan biner) Untuk $p_0 = p_2 = \frac12$
        (kritis), pecahkan persamaan fungsionalnya:
        \[
        H(t) = \frac{1 - \sqrt{1 - t^2}}{t},
        \]
        lalu uraikan dengan \cref{ex:b2:powerseries:catalan} untuk
        memperoleh
        \[
        \P(Y = 2k + 1) = \frac{C_k}{2^{2k+1}},
        \qquad C_k = \frac1{k+1}\binom{2k}k ;
        \]
        periksalah nilai $\P(Y = 1) = \frac12$ dan $\P(Y = 3)
        = \frac18$ lewat pencacahan langsung.
  \item Dengan menurunkan persamaan fungsionalnya di $1^-$,
        tunjukkan bahwa $\E(Y) = \frac{1}{1-m}$ untuk $m < 1$, sedangkan
        kekritisannya memaksa $\E(Y) = \infty$: total keturunan
        kritisnya berhingga hampir pasti dengan rerata tak
        hingga.
  \item Dengan asimtotik binomial pusatnya
        (\cref{ex:b2:comparison:centralbinomial}), tunjukkan
        \[
        \P(Y = 2k+1) \sim \frac{1}{2\sqrt\pi\,k^{3/2}},
        \]
        yakni ekor berat $k^{-3/2}$, lalu simpulkan $\P(Y > n)
        \asymp n^{-1/2}$ (batas atas dan bawah berorde ini
        sudah memadai).
  \item Bandingkan dengan \omterm{pb:b2:proba:1}{jalan acak} setimbang (soal akhir
        pekan pada \cref{ch:b2:proba}): waktu pulang yang pasti tetapi
        dengan rerata tak hingga di sana, total keturunan yang pasti tetapi
        dengan rerata tak hingga di sini, keduanya dengan
        \omterm{def:b2:randomvar:law}{hukum} lokal $n^{-3/2}$. Satu paragraf tentang mengapa
        kekritisannya menghasilkan tanda tangan ini.
\end{enumerate}

\medskip
\noindent\textbf{Bagian IV --- Taksiran Kolmogorov pada
kekritisannya.} Andaikan $m = 1$, $0 < \sigma^2 = G''(1) <
\infty$.
\begin{enumerate}[resume]
  \item Tunjukkan bahwa $G''$ meluas secara \omterm{def:b2:metric:continuity}{kontinu} ke $\intcc01$
        \emph{(taknegatif dan naik dengan limit yang berhingga)} lalu
        simpulkan uraian Taylornya di $1$:
        \[
        G(t) = t + b\,(1-t)^2 + o\bigl((1-t)^2\bigr),
        \qquad b = \frac{G''(1)}2 = \frac{\sigma^2}2 .
        \]
  \item Untuk $t \in \intco01$, tetapkan $h(t) = \dfrac1{1 - G(t)} -
        \dfrac1{1 - t}$. Tunjukkan
        \[
        h(t) = \frac{G(t) - t}{(1 - G(t))(1 - t)}
        \xrightarrow[t\to1^-]{} b .
        \]
  \item Teleskopkan sepanjang iterasi $q_{j+1} = G(q_j)$:
        \[
        \frac1{1 - q_n} = 1 + \sum_{j=0}^{n-1}h(q_j),
        \]
        lalu rampungkan dengan hujah Cesàro bahwa
        \[
        \P(Z_n > 0) = 1 - q_n \sim \frac{2}{\sigma^2\,n}
        \]
        --- yakni \emph{taksiran Kolmogorov}: setiap \omterm{pb:b2:genfun:1}{proses percabangan}
        kritis mati dengan laju semesta $1/n$,
        dan hanya konstantanya yang mengingat \omterm{def:b2:randomvar:law}{hukum} keturunannya.
  \item Periksalah taksiran itu terhadap kasus geometrik
        kritis pada pertanyaan 10.
  \item Simpulkan $\E(Z_n \mid Z_n > 0) = \dfrac{1}{1 - q_n}
        \sim \dfrac{\sigma^2 n}{2}$ \emph{(catat $\E(Z_n
        \mathbf 1_{Z_n>0}) = \E(Z_n) = 1$)}, lalu periksalah
        terhadap pertanyaan 11: dengan disyaratkan pada kesintasannya,
        populasinya tumbuh \emph{secara linear} --- itulah tali
        kritis antara kematian dan ledakan.
\end{enumerate}

\medskip
\noindent\textbf{Bagian V --- Terapan dan sintesisnya.}
\begin{enumerate}[resume]
  \item (Wabah, reaksi berantai) Untuk \omterm{def:b2:randomvar:law}{hukum} keturunan Poisson
        $\mathcal P(\lambda)$ --- setiap kasus menulari
        $\mathcal P(\lambda)$ kasus baru --- tuliskan
        persamaan kepunahannya $q = \eu^{\lambda(q-1)}$ lalu
        pecahkan secara numerik untuk $\lambda = 1.5$ ($q \approx
        0.417$) dan $\lambda = 2$ ($q \approx 0.203$):
        bermula dari satu kasus, wabah besar itu
        \emph{tidak} pasti bahkan ketika $\lambda > 1$. Jelaskan
        mengapa iterasi $q_{n+1} = \eu^{\lambda(q_n - 1)}$
        dari $q_0 = 0$ konvergen ke akar yang benar.
  \item Bermula dari $k$ leluhur alih-alih satu, tunjukkan bahwa
        peluang kepunahannya adalah $q^k$. Terapannya:
        dengan $\lambda = 1.5$, berapa banyak kasus awal yang membuat
        wabah berkemungkinan sedikitnya $99\%$?
  \item (Menyaratkan proses superkritis pada kepunahannya)
        Untuk $m > 1$ dengan peluang kepunahan $q \in
        \intoo01$: buktikan dahulu lewat kecembungannya bahwa $G'(q) < 1$
        di titik tetap terkecilnya, lalu simpulkan
        $q_{\mathrm{ext}} - q_n = O\bigl(G'(q)^n\bigr)$
        (kekonvergenan geometrik, sebagaimana dicontohkan pertanyaan 9).
        Lalu tunjukkan bahwa $\widehat G(t) = G(qt)/q$ merupakan \omterm{ex:b2:powerseries:fibonacci}{fungsi
        pembangkit} sebuah \omterm{def:b2:randomvar:law}{hukum} keturunan yang sahih, dengan rerata $\widehat m =
        G'(q) < 1$: yakni proses pendamping yang subkritis. Sahkanlah
        pada keluarga geometriknya: menyaratkan proses
        $(p, q)$ yang superkritis pada kepunahannya menukar
        $p$ dan $q$. (Pernyataan \omterm{def:b2:metric:complete}{lengkapnya} --- bahwa proses yang disyaratkan itu
        \emph{adalah} proses pendampingnya --- dibuktikan
        pada jilid Tahun ke-3; di sini yang telah disahkan barulah
        bayangan \omterm{ex:b2:powerseries:fibonacci}{fungsi pembangkitnya}.)
  \item Sintesis: susunlah tabel trikotominya --- untuk $m <
        1$, $m = 1$, $m > 1$: nilai $q$; laju $\P(Z_n >
        0)$ atau laju $q - q_n$; $\E(Y)$; ukuran sebuah generasi
        yang sintas. Nyatakan dalam satu kalimat per perkakas bagaimana
        komposisi \omterm{ex:b2:powerseries:fibonacci}{fungsi pembangkit}, kecembungan, Taylor di $1^-$, dan
        perata-rataan Cesàro membawa seluruh soal ini, dan apa
        yang ditambahkan jilid Tahun ke-3 (martingal $Z_n/m^n$ dan
        \omterm{def:b2:randomvar:law}{hukum} limit eksponensial Yaglom).

\end{enumerate}
\end{problem}
